\documentclass[10pt,a4paper]{amsart}
\usepackage[margin=19mm]{geometry}
\usepackage{amsmath,amssymb,mathtools}
\usepackage[scr=rsfs,cal=euler]{mathalfa}
\usepackage{xfrac}
\usepackage[unicode,hidelinks,bookmarks=false]{hyperref}
\newcommand{\ie}{i.e.\ }
\newcommand{\cf}{cf.\ }
\newcommand{\resp}{respectively\ }
\newcommand{\Subsection}{\subsection}
\providecommand{\nobreakdash}{\nobreak\hbox{-}}
\setlength{\emergencystretch}{2em}
\multlinegap0pt
\usepackage{bm}
\usepackage{bbm}
\usepackage{dsfont}
\usepackage{xspace}
\usepackage{cancel}
\usepackage{mathtools}
\usepackage{amsthm}
\makeatletter
\@ifundefined{theorem}{\newtheorem{theorem}{Theorem}}{}
\@ifundefined{claim}{\newtheorem{claim}{Claim}}{}
\@ifundefined{lemma}{\newtheorem{lemma}[theorem]{Lemma}}{}
\@ifundefined{proposition}{\newtheorem{proposition}[theorem]{Proposition}}{}
\makeatother
\DeclareMathOperator{\Isom}{Isom}
\DeclareMathOperator{\comp}{comp}
\title[Uncountably many homogeneous real trees with the same valenc]{Uncountably many homogeneous real trees with the same valence}
\author{P\'en\'elope Azuelos}
\date{Source: arXiv:2511.03722v1 (CC BY 4.0)}
\begin{document}
\maketitle

\noindent\emph{Source and licence.} Uncountably many homogeneous real trees with the same valence. Version arXiv:2511.03722v1 (CC BY 4.0), distributed under the Creative Commons Attribution 4.0 International licence, as recorded in the arXiv licence metadata for this exact version and shown on the abstract page https://arxiv.org/abs/2511.03722v1. Full rights notice: licenses/arxiv-2511.03722-CC-BY-4.0.txt.\par\smallskip

\noindent\textit{Editorial scope.} Counted unit: the Proposition whose source label is \texttt{lem: no isometry} in the article (source0.tex lines 314--320, source label \texttt{lem: no isometry}), delivered together with the article's own complete original proof of it (source0.tex lines 321--391, one \texttt{proof} environment of 71 lines). No mathematics is written, repaired or re-derived: statement and proof are the article's own text line for line. Required context retained uncounted: lem: complexity of the limit (lines 298--300); lem: order preserving isom (lines 307--309). Every definition, display and label of those lines is retained unchanged. Omitted from this package and not redistributed: 1--297, 301--306, 310--313, 392--404. Nothing inside a retained excerpt is omitted or altered: every retained body is delivered exactly as the article's lines are written, so no omission receipt of any kind accompanies this package. Citations: the retained excerpts carry no citation command at all, so no bibliography entry of the article is transcribed and no reference is redistributed. Reference adaptation: none was needed and none was made; every reference inside the delivered unit resolves inside it, so no unresolved reference or citation is produced. Printed slips kept literal: no printed word, symbol, hypothesis or number of the counted statement or of its proof was altered. Layout and class: the article's own class is \texttt{article} with packages including geometry, amsmath, amssymb, amsthm, amsfonts, xcolor, tikz, mathtools, pifont, tabularx, comment, graphicx, hyperref; the wrapper is the lane's pinned \texttt{amsart} editorial wrapper at 10pt on A4, which restates the theorem-like environments the retained text needs and adds the article's own macro definitions that the retained text needs (\texttt{\textbackslash Isom}, \texttt{\textbackslash comp}), with extra wrapper packages (bm, bbm, dsfont, xspace, cancel, mathtools). The delivered numbering is the unit's own. Assets: no retained span contains a figure environment, a graphics-inclusion command, a PDF-page inclusion, a table or a TikZ picture; no image file is copied into this package, the package declares no asset dependency and no image file is redistributed with it. Licence: version arXiv:2511.03722v1 of this article is distributed under the Creative Commons Attribution 4.0 International licence, as recorded in the arXiv licence metadata for this exact version and shown on the abstract page https://arxiv.org/abs/2511.03722v1. A full rights notice is delivered at licenses/arxiv-2511.03722-CC-BY-4.0.txt. Characters: every retained excerpt is pure ASCII, so the delivered engine, XeLaTeX with its default Latin Modern text font, reports no missing character.
\begin{lemma} \label{lem: complexity of the limit}
    Let $(\beta_n)_{n \in \mathbb{N}}$ be a monotone increasing (possibly constant) sequence of countable successor ordinals such that $\beta_1 \geq 1$. Let $\beta \coloneqq \sup_{n \in \mathbb{N}} \beta_n$. Suppose $(a_n)_{n \in \mathbb{N}} \subseteq T_\kappa$ is a sequence such that $a_n \preceq a_{n+1}$ and $\comp(a_n, a_{n+1}) = \beta_n$ for each $n$, and there exists $a \in T_\kappa$ such that $(a_n)_{n \in \mathbb{N}}$ converges to $a$. Then $\comp(a_1,a) = \beta + 1$. Moreover $\comp(a) = \max\{\comp(a_1), \beta + 1\}$.
\end{lemma}

\begin{lemma} \label{lem: order preserving isom}
    Let $\varphi: \mathbb{R} \rightarrow T_\kappa^{[\alpha]}$ be a geodesic line. Then there exists an isometry $f \in \Isom(T_\kappa^{[\alpha]})$ such that $f(\varphi) = L_0$ and $f \circ \varphi((-\infty,0)) = \{c_r : r < 0\}$.
\end{lemma}

\begin{proposition} \label{lem: no isometry}
    Let $\kappa \geq 3$ be a cardinal and let $\alpha_1, \alpha_2 \geq 1$ be countable ordinals such that $\alpha_1 < \alpha_2$.
    \begin{itemize} 
    	\item The $\mathbb{R}$-tree $T_\kappa^{[\alpha_1]}$ is incomplete.
    	\item If $\psi: T_\kappa^{[\alpha_1]} \hookrightarrow T_\kappa^{[\alpha_2]}$ is an isometric embedding, then $\psi$ is not surjective.
    \end{itemize}
\end{proposition}

\begin{proof}
	By Lemma~\ref{lem: order preserving isom}, we can (and do) assume that $\psi$ restricts to the identity on $L_0$. This implies that $\psi$ is order preserving.

    \setcounter{claim}{0}
    \begin{claim}
     \label{claim: induction hyp}
    	Let $\beta$ be a successor ordinal such that $1 \leq \beta \leq \alpha_1$, let $x \in T_\kappa^{[\alpha_1]}$ and let $y \coloneqq \psi(x)$. Then, for all $r > 0$, there exists $a \in T_\kappa^{[\alpha_1]}$ such that:
    	\begin{itemize}
    		\item[i.] $x \prec a$;
    		\item[ii.] $\rho_a \leq \rho_x + r$;
    		\item[iii.] $\comp(x,a) = \beta$;
    		\item[iv.] $\comp(y, \psi(a)) \leq \beta$. 
    	\end{itemize}
    \end{claim}
    \begin{proof}
    	\renewcommand{\qedsymbol}{$\blacksquare$}
    	We proceed by transfinite induction on $\beta$. Suppose $\beta = 1$ and let $c \in C_\kappa - \{0\}$. If there exists $\varepsilon > 0$ such that $x((\rho_{x} - \varepsilon, \rho_{x})) = \{0\}$ then let $a': (-\infty, \rho_{x} + 1) \rightarrow C_\kappa$ be defined by 
    	\[
    	a'(t) =
    	\begin{cases}
    		x(t) \quad &\text{if } t < \rho_{x};\\
    		c &\text{otherwise.}
    	\end{cases}
    	\]
    	Otherwise, define $a': (-\infty, \rho_{x} + 1) \rightarrow C_\kappa$ by
    	\[ 
    	a'(t) =
    	\begin{cases}
    		x(t) \quad &\text{if } t < \rho_{x}; \\
    		0 &\text{otherwise.}
    	\end{cases}
    	\]
    	In either case $P_{a'} = P_{x} \cup \{\rho_{x}\}$. Moreover, if $a \in T_\kappa$ is such that $x \prec a \preceq a'$, then $P_{a} \cap [\rho_{x}, \rho_{a}] =\{\rho_{x}\}$, so $\comp(x, a) = 1$.
    	Let $r' > 0$ be such that $r' \leq r$ and, if $P_{\psi(a')} \cap (\rho_{y}, \rho_{\psi(a')}) \neq \emptyset$, then $\rho_{y} + r' \leq \min P_{\psi(a')} \cap (\rho_{y}, \rho_{\psi(a')})$.
    	Let $b \coloneqq \psi(a')|_{(-\infty, \rho_{y} + r')}$. Then $P_b \subseteq P_{y} \cup \{\rho_{y}\}$ so $\comp(y,b) \leq 1$.
    	Since $\psi\big(T_\kappa^{[\alpha_1]}\big)$ is connected, there exists $a \in T_\kappa^{[\alpha_1]}$ be such that $b = \psi(a)$. Then $x \prec a \preceq a'$ so $\comp(x, a) = 1$ and $\rho_{a} = \rho_{x} + d(x, a) = \rho_{x} + r' \leq \rho_{x} + r$.
    	
    	Now suppose $\beta = \gamma + 1> 1$ and the claim holds for all successor ordinals $1 \leq \beta' < \beta$. If $\gamma$ is a successor ordinal then let $\gamma_n \coloneqq \gamma$ for all $n \in \mathbb{N}$. If $\gamma$ is a limit ordinal, then let $(\gamma_n)_{n \in \mathbb{N}}$ be a strictly increasing sequence of successor ordinals such that $\sup_{n \in \mathbb{N}} \gamma_n = \gamma$.  We next define sequences $(a_n)_{n \in \mathbb{N}} \subseteq T_\kappa^{[\alpha_1]}$ and $(b_n)_{n \in \mathbb{N}} \subseteq T_\kappa^{[\alpha_2]}$ recursively such that, for all $n \in \mathbb{N}$:
    	\begin{enumerate}
    		\item $\psi(a_n) = b_n$;
    		\item $x \preceq a_n \prec a_{n+1}$ and $y \preceq b_n \prec b_{n+1}$;
    		\item $\comp(a_n,a_{n+1}) = \gamma_{n}$ and $\comp(b_n,b_{n+1}) \leq \gamma_{n}$;
    		\item $\rho_{a_n} \leq \rho_{x} + \sum_{i=1}^n r/2^i$ and $\rho_{b_n} \leq \rho_{y} + \sum_{i=1}^n r/2^i$.
    	\end{enumerate}
    	Let $a_1 \coloneqq x$ and $b_1 \coloneqq y$. Items 1,2 and 4 are satisfied. Let $n \in \mathbb{N}$ and suppose that $a_1, \dots, a_n \in T_\kappa^{[\alpha_1]}$ and $b_1, \dots, b_n \in T_\kappa^{[\alpha_2]}$ have been defined. Applying the induction hypothesis to $(\gamma_{n}, a_n, r/2^{n+1})$, there exists $a_{n+1} \in T_\kappa^{[\alpha_1]}$ such that 
    	\begin{itemize} 
    		\item[i.] $a_n \prec a_{n+1}$;
    		\item[ii.] $\rho_{a_{n+1}} \leq \rho_{a_n} + r/2^{n+1} \leq \rho_{x_1} + \sum_{i=1}^{n+1} r/2^i$;
    		\item[iii.] $\comp(a_n, a_{n+1}) = \gamma_{n}$;
    		\item[iv.] if $b_{n+1} \coloneqq \psi(a_{n+1})$, then $\comp(b_n, b_{n+1}) \leq \gamma_{n}$.
    	\end{itemize}
    	Thus Items 1-3 are satisfied. Moreover 
    	\[
    	\rho_{b_{n+1}} = \rho_{b_n} + d(b_n, b_{n+1}) \leq \rho_{y} + \sum_{i=1}^n \frac{r}{2^i} + d(a_n, a_{n+1}) \leq \rho_{y} + \sum_{i=1}^{n+1} \frac{r}{2^i},
    	\]
    	so Item~4 holds.
    	
    	It follows from Items 2 and 4 that the sequences $(a_n)_{n \in \mathbb{N}}$ and $(b_n)_{n \in \mathbb{N}}$ are strictly increasing and Cauchy. Let $a \coloneqq \lim_{n \rightarrow \infty} a_n \in T_\kappa$, $b \coloneqq \lim_{n \rightarrow \infty} b_n \in T_\kappa$. By Lemma~\ref{lem: complexity of the limit} and Item 3, $\comp(x,a) = \gamma + 1 = \beta$, $\comp(a) \leq \alpha_1$, $\comp(y,b) \leq \gamma + 1 < \alpha_2$ and $\comp(b) \leq \alpha_2$. Since $\psi$ is continuous, $\psi(a) = b$.
    	Moreover $x \prec a$ and $\rho_{a} = \lim_{n \rightarrow \infty} \rho_{a_n} \leq \rho_{x} + r$.
    \end{proof}

    Let $(\beta_n)_{n \in \mathbb{N}}$ be a monotone increasing (possibly constant) sequence of successor ordinals such that $\beta_n \geq 1$ for each $n$ and $\sup_{n \in \mathbb{N}} \beta_n = \alpha_1$.
    It follows from Claim~\ref{claim: induction hyp} (using the same argument as in the construction of $(a_n)_{n \in \mathbb{N}}$ and $(b_n)_{n \in \mathbb{N}}$ in the proof of that claim) that there exist sequences $(a_n)_{n \in \mathbb{N}} \subseteq T_\kappa^{[\alpha_1]}$ and $(b_n)_{n \in \mathbb{N}} \subseteq T_\kappa^{[\alpha_2]}$ such that $a_1 = b_1 = c_0$ and, for each $n \in \mathbb{N}$:
    \begin{enumerate}
            \item $\psi(a_n) = b_n$;
            \item $a_n \prec a_{n+1}$ and $b_n \prec b_{n+1}$;
            \item $\comp(a_n,a_{n+1}) = \beta_{n}$ and $\comp(b_n, b_{n+1}) \leq \beta_{n}$;
            \item $\rho_{a_n} \leq \rho_{a_1} + \sum_{i=1}^n r/2^i$ and $\rho_{b_n} \leq \rho_{b_1} + \sum_{i=1}^n r/2^i$.
    \end{enumerate}
    Items 2 and 4 imply that $(a_n)_{n \in \mathbb{N}}$ and $(b_{n \in \mathbb{N}})_{n \in \mathbb{N}}$ are Cauchy and Lemma~\ref{lem: complexity of the limit} implies that the complexity of $a \coloneqq \lim_{n \in \mathbb{N}} a_n \in T_\kappa$ is $\alpha_1 + 1$ and the complexity of $b \coloneqq \lim_{n \rightarrow \infty} b_n \in T_\kappa$ is $\leq \alpha_1 + 1 \leq \alpha_2$. Therefore $a \notin T_\kappa^{[\alpha_1]}$, so $T_\kappa^{[\alpha_1]}$ is incomplete, and $b = \lim_{n \rightarrow \infty} \psi(a_n) \in T_\kappa^{[\alpha_2]} - \psi(T_\kappa^{[\alpha_1]})$, so $\psi$ is not surjective.
\end{proof}

\end{document}
